\section{Discussion}

The executor result supports a narrow systems observation: contiguous packed
blocks can outperform optimized dense execution once enough blocks are
removed. It rejects the stronger project goal for the tested setup because the
models stop meeting the quality gate much earlier. The direction did not drift;
it progressively isolated why theoretical sparsity did not become useful
latency.

Several attempted remedies did not close the gap. Calibration-only LSH and
hot/cold reorderings improved some block retention or perplexity but did not
make 20--30\% Llama conditions pass. Direct NEON and coalesced Accelerate
variants were slower than the existing native B8 implementation. Custom INT8
experiments won against weak matched INT8 controls but not against the strong
FP32 Accelerate latency, while a minimal llama.cpp Q8$_0$ B32 probe did not
reliably cross dense parity through 80\%. These negative controls discourage
presenting backend tuning as the missing scientific contribution.

\subsection{Threats to validity}

The quality screen is small and uses perplexity rather than downstream task
accuracy. Oracle contribution scoring reads information unavailable to a
deployable predictor. Runtime masks are favorable and synthetic. Measurements
cover one CPU, one thread, one token, FP32 kernel arithmetic, and only two FFN
geometries. Reported ranges are repeated-run observations, not confidence
intervals. These constraints limit generalization but make the bounded
conclusion auditable.

\subsection{Next research question}

The evidence shifts the bottleneck from execution toward representation. A
credible continuation must move the quality frontier from 10--20\% toward the
40--50\% mechanical break-even. Training-aware block activation, joint
selector/layout learning, or architectures with native grouped FFNs are more
plausible than another local kernel micro-optimization. Any continuation should
pre-register a stop condition and compare with modern dense and sparse ARM
baselines.
